\chapter{Melioidosis}

\section*{Definition and Overview}
Melioidosis, also known as Whitmore's disease, is a severe and potentially life-threatening infectious disease caused by the Gram-negative, facultative intracellular bacterium \textit{Burkholderia pseudomallei}. Found naturally in soil and surface water, the organism is endemic to tropical and subtropical regions. The disease is highly significant clinically due to its diverse and often deceptive presentations, ranging from localized cutaneous abscesses to fulminant, rapidly progressive septicaemia with multi-organ failure. It behaves as an opportunistic pathogen, frequently affecting individuals with underlying chronic conditions, and represents a major public health concern in endemic zones, particularly in northern Australia and Southeast Asia.

\section*{Epidemiology}
Melioidosis is highly endemic in northern Australia (extending across the Top End of the Northern Territory, the Kimberley region of Western Australia, and northern Queensland) and throughout Southeast Asia (notably Thailand, Singapore, and Malaysia). It is also increasingly recognized in parts of South Asia, Central and South America, and Africa. In Australia, the disease exhibits a marked seasonal pattern, with the vast majority of cases occurring during the wet season (from November to April) when rainfall, high humidity, and tropical storms disturb the soil and bring the bacteria to the surface. It can affect individuals of any age, but the peak incidence is in adults aged between 40 and 60 years, with a higher prevalence in males, largely due to occupational and recreational exposure to soil and mud. Indigenous Australian populations are disproportionately affected, primarily due to higher rates of predisposing comorbidities and environmental exposures.

\section*{Aetiology and Pathophysiology}
The causative agent, \textit{Burkholderia pseudomallei}, is an environmental saprophyte that persists in wet soils and surface agricultural waters.
\begin{itemize}
    \item \textbf{Modes of transmission:} Infection occurs primarily through direct inoculation of contaminated soil or water through skin abrasions, cuts, or puncture wounds. It can also be acquired via inhalation of aerosolised bacteria during severe weather events (such as cyclones or monsoons) or, less commonly, through ingestion of contaminated untreated water.
    \item \textbf{Pathogenic mechanisms:} Once inside the host, \textit{B. pseudomallei} is capable of surviving and replicating within both phagocytic and non-phagocytic cells. It utilises a capsule to resist phagocytosis and employs type III secretion systems to escape endosomes, moving intercellularly by inducing cell-to-cell fusion, which results in the formation of multinucleated giant cells.
    \item \textbf{Host risk factors:} The most critical predisposing factor is \textbf{diabetes mellitus}, which impairs neutrophil function and intracellular killing capacity, multiplying the risk of contracting the disease. Other significant risk factors include \textbf{hazardous alcohol consumption}, \textbf{chronic kidney disease}, \textbf{chronic lung diseases} (such as chronic obstructive pulmonary disease), \textbf{thalassaemia}, and states of \textbf{immunosuppression} (e.g., oncology treatment or high-dose corticosteroid therapy).
\end{itemize}

\section*{Clinical Features}
The clinical presentation of melioidosis is highly variable and depends on the route of infection, bacterial load, and host immune status.
\begin{itemize}
    \item \textbf{Acute septicaemia:} Typically presents with sudden onset of high fever, rigors, severe headache, confusion, muscle aches, and rapid progression to septic shock with hypotension.
    \item \textbf{Pulmonary infection:} This is the most common presentation, ranging from a mild bronchitis-like illness to severe, necrotising pneumonia. Symptoms include a productive or non-productive cough, high fever, pleuritic chest pain, and progressive shortness of breath.
    \item \textbf{Localised abscesses:} Cutaneous ulcers, non-healing wounds, and subcutaneous abscesses can develop at the site of inoculation. Visceral abscesses are common, particularly in the spleen, liver, and prostate, often presenting with localised pain and persistent fever.
    \item \textbf{Parotid abscesses:} A classic clinical sign in paediatric patients, particularly in Southeast Asia, presenting as unilateral or bilateral painful swelling of the parotid gland.
    \item \textbf{Neurological manifestations:} Known as neuromelioidosis, this presentation can include brainstem encephalitis, meningitis, or flaccid paralysis, which carry a particularly high morbidity.
\end{itemize}

\section*{Differential Diagnosis}
Because of its wide range of symptoms, melioidosis is known as the ``remarkable imitator'' and must be distinguished from several other conditions:
\begin{itemize}
    \item \textbf{Tuberculosis:} Shares features of chronic cough, fever, night sweats, weight loss, and cavitary lung lesions on chest imaging.
    \item \textbf{Bacterial pneumonia:} Community-acquired pneumonias caused by pathogens such as \textit{Streptococcus pneumoniae}, \textit{Klebsiella pneumoniae}, or \textit{Legionella} species.
    \item \textbf{Systemic sepsis:} Sepsis caused by common pathogens like \textit{Staphylococcus aureus} or other Gram-negative bacilli (e.g., \textit{Escherichia coli}, \textit{Pseudomonas aeruginosa}).
    \item \textbf{Leptospirosis and typhus:} Other tropical febrile illnesses with similar geographic distributions and risk profiles involving exposure to mud and fresh water.
    \item \textbf{Fungal infections:} Deep systemic mycoses that present with chronic nodular or cavitary pulmonary disease and abscess formation.
\end{itemize}

\section*{When to Seek Medical Care}
Individuals living in or visiting tropical regions where melioidosis is endemic should seek prompt medical evaluation if they develop a high fever, a cough that does not resolve, or painful, swollen skin lesions, especially if they have predisposing risk factors such as diabetes or kidney disease.

\textbf{Call 000 (or local emergency services) for:}
\begin{itemize}
    \item Severe shortness of breath or difficulty breathing.
    \item High fever accompanied by confusion, severe lethargy, or loss of consciousness.
    \item Signs of septic shock, including a rapid heart rate, low blood pressure, cold, clammy, or pale skin, and rapid breathing.
    \item Severe chest pain.
\end{itemize}

\section*{Diagnosis and Investigations}
Early and accurate diagnosis is critical, requiring a high index of clinical suspicion in endemic areas.
\begin{itemize}
    \item \textbf{Microbiological culture:} The gold standard of diagnosis. \textit{Burkholderia pseudomallei} is isolated from clinical specimens including blood, sputum, urine, abscess aspirates, throat swabs, or skin lesions. Laboratory staff must be warned if melioidosis is suspected due to the biosafety risk of handling the bacterium.
    \item \textbf{Imaging studies:} A chest X-ray or CT scan is essential to evaluate for pulmonary consolidation, nodules, or cavitary lesions. Abdominal ultrasound or CT imaging should be performed in all confirmed cases to screen for clinically silent abscesses in the liver, spleen, and prostate.
    \item \textbf{Molecular testing:} Polymerase chain reaction (PCR) assays can facilitate rapid detection of \textit{B. pseudomallei} DNA directly from blood or clinical samples, providing results much faster than culture.
    \item \textbf{Serological testing:} The indirect haemagglutination assay (IHA) can detect antibodies, but its utility is limited in highly endemic areas where a significant proportion of the healthy population has background seropositivity.
\end{itemize}

\section*{Management}
Treatment of melioidosis is prolonged and divided into two distinct phases to prevent relapse.
\begin{itemize}
    \item \textbf{Intensive phase (Intravenous):} Initiated immediately with intravenous ceftazidime or meropenem for a minimum of 10 to 14 days. For severe cases, such as those with neuromelioidosis, osteomyelitis, or deep-seated visceral abscesses, the intravenous phase is extended to 4 weeks or longer.
    \item \textbf{Eradication phase (Oral):} Follows the intensive phase to clear residual intracellular bacteria. The standard regimen is oral trimethoprim-sulfamethoxazole (co-trimoxazole) taken for 3 to 6 months. In cases where this is contraindicated, co-amoxiclav (amoxicillin-clavulanic acid) may be used.
    \item \textbf{Surgical intervention:} Drainage of large abscesses (e.g., prostatic or large joint collections) may be required to facilitate clinical resolution.
    \item \textbf{Supportive care:} Critical care management in an intensive care unit (ICU) is required for patients presenting with septic shock, including vasoactive support and invasive mechanical ventilation.
\end{itemize}

\section*{Complications}
Despite treatment, the clinical course of melioidosis can be complicated by severe pathology.
\begin{itemize}
    \item Septic shock and progressive multi-organ dysfunction syndrome (MODS).
    \item Relapse of the infection, which occurs in up to 10\% of cases, particularly in patients who do not complete the long eradication course of oral antibiotics.
    \item Chronic pulmonary dysfunction, including bronchiectasis, secondary to severe cavitary pneumonia.
    \item Permanent neurological deficits (such as cranial nerve palsies or limb weakness) following neuromelioidosis.
    \item Secondary bacterial infections during prolonged hospitalisation.
\end{itemize}

\section*{Prognosis and Outcomes}
The prognosis for melioidosis depends heavily on the patient's baseline risk factors, the severity of presentation, and the promptness of initiating appropriate antibiotic therapy. In Australia, where access to advanced intensive care is widely available, the overall mortality rate is approximately 10\% to 15\%. However, in resource-limited settings across Southeast Asia, mortality rates can exceed 40\%. Patients presenting in septic shock have a poor prognosis, with mortality rates remaining high even with optimal care. Relapse rates are low (under 5\%) when the full course of intravenous and oral eradication therapy is adhered to, but rise significantly if treatment is truncated.

\section*{Prevention and Self-Care}
There is currently no vaccine available for melioidosis. Prevention relies on reducing environmental exposure.
\begin{itemize}
    \item Avoid direct contact with soil, mud, and surface water during and immediately after the wet season, especially for individuals with diabetes, kidney disease, or open wounds.
    \item Wear protective clothing, including closed waterproof footwear (such as boots) and gloves, when gardening, farming, or engaging in outdoor activities.
    \item Clean and cover any cuts, scratches, or skin abrasions immediately with waterproof dressings if exposure to soil or water has occurred.
    \item Drink only treated, safe drinking water, and avoid using untreated bore water for drinking or personal hygiene.
    \item Ensure optimal management of predisposing conditions, particularly maintaining strict glycaemic control in diabetes.
\end{itemize}

\section*{Sources and Further Reading}
The following reputable organisations provide further information and clinical guidelines on melioidosis:
\begin{itemize}
    \item Northern Territory Government Department of Health. \textit{Melioidosis clinical guidelines and public health resources.} https://health.nt.gov.au/
    \item Healthdirect Australia. \textit{Consumer health information and advice on melioidosis.} https://www.healthdirect.gov.au/
    \item Centers for Disease Control and Prevention (CDC). \textit{Detailed disease overview, transmission, and prevention information.} https://www.cdc.gov/
    \item World Health Organization (WHO). \textit{Global epidemiology and resources on tropical diseases.} https://www.who.int/
    \item MSD Manual (Professional Edition). \textit{Clinical manual detailing the pathophysiology, diagnosis, and treatment of melioidosis.} https://www.msdmanuals.com/
\end{itemize}
